\section{\MLIPE{}: Domain-Specific Extensions for \Concon{}}
\label{sec:mlipe}
%% ====================================================================

The \MLIPsOntology{}, SHACL shapes, and \MLIPsKG{} are hosted on \Concon{}, a
collaborative knowledge graph platform that provides general-purpose
infrastructure for team-based management of RDF knowledge
graphs~\cite{hernandez2026concon}. \Concon{} handles authentication,
hierarchical team structures, role-based access control, per-dataset SPARQL
endpoints, SHACL-driven entity editing, literature management, and AI-assisted
knowledge extraction. These features are domain-agnostic and serve
\Concon{}'s other deployments in computational design and construction,
circular manufacturing, and university teaching.

\MLIPE{} (Machine Learned Interatomic Potentials Environment) is a set of
extension modules for \Concon{} that add features specific to MLIP and
atomistic simulation research. The distinction is important: features that
benefit any knowledge graph user belong in \Concon{}'s core, while features
tied to the MLIP domain belong in \MLIPE{}.

\subsection{\MLIPE{}-Specific Modules}

The following modules are specific to \MLIPE{} and not part of
\Concon{}'s core:

\begin{enumerate}
\item \emph{Custom SPARQL extension functions.} Domain-specific computations
  that operate on atomistic data stored in the knowledge graph: radial
  distribution functions, phonon spectra calculations, and energy/force
  analysis from stored atomic configurations. These are implemented as
  Oxigraph custom functions and exposed through the SPARQL query interface.

\item \emph{Simulation data handling.} Support for multi-level storage of
  simulation data: meta-level descriptions of MLIP algorithms and
  hyperparameters alongside simulation-level per-atom states (positions,
  forces, velocities) and system-level properties (total energy, stress
  tensors, lattice vectors). All numeric values are stored as typed RDF
  literals with units annotated via QUDT~\cite{hodgson2014qudt}.

\item \emph{\SemASE{} integration.} A dedicated API for ingesting
  simulation data from \SemASE{} (Section~\ref{sec:semase}), including
  batch import of atomic configurations from ASE-based workflows.

\item \emph{MLIP-specialized AI extraction agent.} A domain-specialized
  agent that reads papers from the literature database and produces structured
  statements about MLIP algorithms, training datasets, hyperparameters, and
  reported accuracy metrics, stored in named graphs with provenance metadata
  for expert curation.

\item \emph{External data caching.} For large simulation datasets hosted in
  external archives (e.g., DaRUS~\cite{darus2024}), \MLIPE{} caches data
  locally so that custom SPARQL functions can operate on it directly from
  the query interface without repeated downloads.
\end{enumerate}

\subsection{Features Provided by \Concon{}'s Core}

The following features, while essential for the MLIP use case, are provided
by \Concon{}'s core and shared with all deployments:

\begin{itemize}
\item Team hierarchy and permissions (META-LEARN work packages)
\item Per-dataset SPARQL endpoints and YASGUI editor
\item SHACL-driven entity editing (for the \MLIPsOntology{} shapes)
\item Literature management with GROBID-based PDF extraction
\item General AI-assisted extraction pipeline
\item File storage and metadata catalogs
\end{itemize}

%% ====================================================================
